\documentclass[11pt,a4paper]{article}

\usepackage[utf8]{inputenc}
\usepackage[T1]{fontenc}
\usepackage[margin=1in]{geometry}
\usepackage{amsmath,amssymb}
\usepackage[protrusion=true,expansion=false]{microtype}
\usepackage{graphicx}
\usepackage{xcolor}
\usepackage{listings}
\usepackage{hyperref}
\usepackage{booktabs}
\usepackage{enumitem}
\usepackage{tcolorbox}
\usepackage{titlesec}

\hypersetup{
  colorlinks=true,
  linkcolor=blue!60!black,
  urlcolor=blue!60!black,
  citecolor=blue!60!black
}

\definecolor{codebg}{rgb}{0.97,0.97,0.97}
\definecolor{codecomment}{rgb}{0.4,0.6,0.4}
\definecolor{codekw}{rgb}{0.0,0.2,0.6}
\definecolor{codestr}{rgb}{0.6,0.2,0.2}

\lstdefinestyle{pystyle}{
  language=Python,
  backgroundcolor=\color{codebg},
  basicstyle=\ttfamily\footnotesize,
  keywordstyle=\color{codekw}\bfseries,
  commentstyle=\color{codecomment}\itshape,
  stringstyle=\color{codestr},
  breaklines=true,
  breakatwhitespace=false,
  showstringspaces=false,
  numbers=left,
  numberstyle=\tiny\color{gray},
  numbersep=5pt,
  frame=single,
  framesep=4pt,
  rulecolor=\color{gray!50},
  xleftmargin=14pt,
  tabsize=4,
}
\lstset{style=pystyle}

\newtcolorbox{srcbox}[1][]{%
  colback=blue!2!white,
  colframe=blue!40!black,
  boxrule=0.4pt,
  arc=2pt,
  left=6pt,right=6pt,top=4pt,bottom=4pt,
  fonttitle=\bfseries\small,
  title=Where to look in the official DINOv3 repository, #1
}

\titleformat{\section}{\Large\bfseries\sffamily}{\thesection}{1em}{}
\titleformat{\subsection}{\large\bfseries\sffamily}{\thesubsection}{1em}{}
\titleformat{\subsubsection}{\normalsize\bfseries\sffamily}{\thesubsubsection}{1em}{}

\title{\textbf{Tailoring DINOv3 for Semantic Segmentation}\\[0.4em]
\large An implementation walkthrough of the linear-head and Mask2Former heads
released in \texttt{facebookresearch/dinov3}\\[0.6em]
\small Coursework report --- DINOv3 study project}
\author{}
\date{}

\begin{document}
\maketitle

\begin{abstract}
\noindent
This report describes how the official DINOv3 codebase
(\url{https://github.com/facebookresearch/dinov3}) adapts a frozen DINOv3
ViT backbone to semantic segmentation. Two adaptations are released by Meta
and are covered here in detail: a minimal \textbf{linear head}
(a single $1{\times}1$ convolution on patch features), trained end-to-end on
ADE20K, and a heavyweight \textbf{Mask2Former (M2F)} head fed by a
\textbf{DINOv3-Adapter} that turns the single-scale ViT into a multi-scale
feature pyramid. The first is fully trainable from the released code; the
second is shipped as inference-only with a pretrained checkpoint that
achieves the headline 55.9 mIoU on ADE20K reported in Meta's model card. The
report is grounded in the source code: every claim cites a file and line
range so a reader can verify it. We do not reproduce Meta's source code in
this document, only describe and reference it.
\end{abstract}

\tableofcontents
\newpage

% =============================================================================
\section{What semantic segmentation needs from a backbone}

Semantic segmentation requires a per-pixel class prediction. A ViT backbone
naturally produces per-\emph{patch} features --- for a $512{\times}512$ input
with patch size 16, a $32{\times}32$ grid of 1024 feature vectors. To get a
per-pixel prediction at resolution $H{\times}W$, the segmentation pipeline
must:
\begin{enumerate}[leftmargin=1.6em,itemsep=0pt]
  \item extract patch features from the backbone, possibly from several
        depths in the network;
  \item combine and project them into class logits per patch;
  \item upsample the patch-level logits to pixel resolution.
\end{enumerate}
The two heads in DINOv3 differ in how much work they do at each step.
The \emph{linear} head is intentionally trivial (concatenate features along
the channel axis, BatchNorm, $1{\times}1$ Conv, bilinear upsample); its
purpose is to demonstrate that the backbone features are already
class-discriminative \emph{without} a heavy decoder. The \emph{M2F} head, by
contrast, treats segmentation as mask classification: it learns a fixed set
of object queries and predicts masks per query, which means it needs a
multi-scale feature pyramid the ViT does not natively provide --- hence the
\emph{DINOv3-Adapter}. We unpack both in turn.

% =============================================================================
\section{Where the segmentation code lives}

\begin{srcbox}
\texttt{dinov3/eval/segmentation/} (the whole subtree).\\
Plus \texttt{dinov3/hub/segmentors.py} for the prebuilt M2F segmentor loader,
and \texttt{dinov3/eval/utils.py:ModelWithIntermediateLayers}.
\end{srcbox}

A summary of the directory:
\begin{description}[leftmargin=1.6em,style=nextline]
  \item[\texttt{models/\_\_init\_\_.py}] the dispatcher
        \texttt{build\_segmentation\_decoder()}: chooses linear vs.\ M2F,
        wires the backbone to the chosen head, and returns a single
        \texttt{FeatureDecoder} module that exposes a uniform
        \texttt{forward()/predict()} interface to the trainer and the
        evaluator.
  \item[\texttt{models/heads/}] the heads themselves --- \texttt{linear\_head.py}
        and \texttt{mask2former\_head.py} (plus its \texttt{pixel\_decoder.py}
        and \texttt{mask2former\_transformer\_decoder.py}).
  \item[\texttt{models/backbone/dinov3\_adapter.py}] the multi-scale adapter
        used only by M2F.
  \item[\texttt{models/utils/}] support layers shared by M2F: deformable
        attention (custom CUDA op), positional encodings, helper transformer
        blocks, batch-norm wrappers.
  \item[\texttt{train.py}, \texttt{eval.py}, \texttt{run.py}] train and eval
        loops; \texttt{run.py} is the entry point both invoke.
  \item[\texttt{inference.py}] the sliding-window inference utility used at
        evaluation time and by our custom-image script.
  \item[\texttt{loss.py}] cross-entropy and Dice losses with class weighting
        and ignore-index support.
  \item[\texttt{metrics.py}] mIoU computation.
  \item[\texttt{transforms.py}] paired image+mask data augmentations.
  \item[\texttt{configs/config-ade20k-linear-training.yaml}] the training
        recipe for the linear head.
  \item[\texttt{configs/config-ade20k-m2f-inference.yaml}] the
        \emph{inference-only} recipe for M2F.
\end{description}

% =============================================================================
\section{The frozen-backbone interface}
\label{sec:backbone-interface}

\begin{srcbox}
\texttt{dinov3/eval/utils.py} (the class \texttt{ModelWithIntermediateLayers}),\\
\texttt{dinov3/models/vision\_transformer.py:get\_intermediate\_layers}.
\end{srcbox}

Both heads consume the same backbone interface: a list of patch-feature maps
extracted from selected transformer blocks. The plumbing is:

\begin{enumerate}[leftmargin=1.6em,itemsep=0pt]
  \item The DINOv3 ViT exposes
        \texttt{get\_intermediate\_layers(x, n=..., reshape=True,
        return\_class\_token=...)} which returns, for each requested block
        index, a tensor of shape $(B, d, H/P, W/P)$ if \texttt{reshape=True}
        (so it is already a 2D feature map), and optionally the corresponding
        $(B, d)$ CLS token.
  \item \texttt{ModelWithIntermediateLayers} wraps that in a callable module
        and adds an autocast context. It also calls
        \texttt{requires\_grad\_(False)} on the wrapped backbone, freezing it
        for the segmentation experiments.
  \item Which blocks to extract is controlled by the
        \texttt{BackboneLayersSet} enum in
        \texttt{eval/segmentation/models/\_\_init\_\_.py}:
        \begin{itemize}[leftmargin=1.6em,itemsep=0pt]
          \item \texttt{LAST}: only the last block. Used by the linear head.
          \item \texttt{FOUR\_LAST}: the last four blocks.
          \item \texttt{FOUR\_EVEN\_INTERVALS}: four blocks evenly spaced.
                For ViT-L/16 (24 blocks) the indices are
                \texttt{[4, 11, 17, 23]} -- this is what the M2F config
                uses.
        \end{itemize}
\end{enumerate}

In short: the head receives a list of one or four 2D feature maps (and
optionally CLS tokens) at the patch-grid resolution. The head's job is to
turn those into class logits at pixel resolution.

% =============================================================================
\section{The linear segmentation head}
\label{sec:linear}

\begin{srcbox}
\texttt{dinov3/eval/segmentation/models/heads/linear\_head.py},\\
\texttt{configs/config-ade20k-linear-training.yaml},\\
\texttt{eval/segmentation/train.py} (training loop),\\
\texttt{eval/segmentation/loss.py}, \texttt{eval/segmentation/metrics.py}.
\end{srcbox}

\subsection{Architecture}

The linear head is approximately 90 lines of code and is essentially a
single $1{\times}1$ convolution. Its construction takes:
\begin{itemize}[leftmargin=1.6em,itemsep=0pt]
  \item \texttt{in\_channels}: a list of per-feature-map dimensions. For
        the default \texttt{LAST} configuration with a ViT-L backbone the
        list is just \texttt{[1024]}; for \texttt{FOUR\_EVEN\_INTERVALS}
        it would be \texttt{[1024, 1024, 1024, 1024]}.
  \item \texttt{n\_output\_channels}: number of classes
        ($150$ for ADE20K, plus an ignore index handled by the loss).
  \item \texttt{use\_batchnorm} (default \texttt{true}): a
        \texttt{nn.SyncBatchNorm} on the concatenated channels.
  \item \texttt{use\_cls\_token} (default \texttt{false} for the ADE20K
        config): if true, the CLS token is broadcast spatially and
        concatenated with the patch tokens, doubling the channel count.
  \item \texttt{dropout} (default $0.1$): \texttt{nn.Dropout2d} applied
        before the $1{\times}1$ convolution.
\end{itemize}

The forward pass is:

\begin{enumerate}[leftmargin=1.6em,itemsep=0pt]
  \item If \texttt{use\_cls\_token}, broadcast each block's CLS token
        $(B,d)\to(B,d,h,w)$ and concatenate with that block's patch map
        along the channel axis.
  \item Bilinearly upsample every block's feature map to the spatial
        resolution of the first block (this is a no-op when only the last
        block is extracted) and concatenate them along the channel axis.
        The result is a single dense map $(B, \sum d_i, h, w)$.
  \item Dropout2d $\to$ BatchNorm $\to$ Conv2d$_{1\times1}$.
  \item For prediction, bilinearly upsample the $(B, K, h, w)$ logits to
        the requested target size (typically the input image's spatial
        resolution).
\end{enumerate}

The intent of this head is pedagogical and methodological: it shows that
DINOv3's patch features are linearly separable into the 150 ADE20K classes
to a useful degree without a heavy decoder. The head adds only $K \cdot
\sum d_i$ parameters (e.g.\ $150 \cdot 1024 = 153{,}600$ for ViT-L/16 LAST),
plus the BatchNorm scale/shift -- against the millions in the backbone.

\subsection{Training recipe}

Reading \texttt{configs/config-ade20k-linear-training.yaml} top-to-bottom
gives the exact recipe:
\begin{itemize}[leftmargin=1.6em,itemsep=0pt]
  \item \textbf{Hardware:} 8 GPUs, batch size 2 each (effective batch 16).
  \item \textbf{Iterations:} 40\,000 total.
  \item \textbf{Optimizer:} AdamW, $\beta_1{=}0.9$, $\beta_2{=}0.999$,
        weight decay $10^{-3}$, learning rate $10^{-3}$. Gradient clipping
        disabled (\texttt{"inf"}).
  \item \textbf{Schedule:} \texttt{WarmupOneCycleLR}: linear warmup over
        1500 iterations from $10^{-9}$ up to peak $10^{-3}$, then a single
        cosine descent. \texttt{pct\_start} is $0$ which makes the
        ``warm-up'' part part of the descent path; the effective peak
        occurs immediately after warm-up.
  \item \textbf{Loss:} Cross-entropy weight $1.0$, Dice weight $0.0$.
        (Dice is implemented in \texttt{loss.py} but disabled in this
        config.)
  \item \textbf{Class handling:} \texttt{reduce\_zero\_label: true} maps
        ADE20K's class 0 (``unlabelled'') to the ignore index, so the
        head learns 150 classes only.
  \item \textbf{Augmentations:} random scale ratio in $[0.5, 2.0]$,
        random crop to $512{\times}512$, random horizontal flip with
        probability $0.5$. The whole transform is implemented as a paired
        operation on (image, mask) so the mask follows every geometric
        change.
  \item \textbf{Eval:} sliding-window inference at $512{\times}512$ with
        stride $341$ ($\approx \tfrac{2}{3}$ overlap), evaluated every
        5000 iterations. No test-time augmentation. Metric: mean IoU
        over the 150 valid classes.
\end{itemize}

\subsection{Sliding-window inference}
\label{ssec:slide}

\begin{srcbox}
\texttt{dinov3/eval/segmentation/inference.py:slide\_inference}.
\end{srcbox}

The eval procedure is sliding-window with overlap. Given an input
$(1, 3, H, W)$, a crop size $h_c{\times}w_c$, and a stride $h_s{\times}w_s$,
the function:
\begin{enumerate}[leftmargin=1.6em,itemsep=0pt]
  \item Computes the grid of crops needed to tile the image, where the
        last column/row is shifted in to avoid going out of bounds.
  \item For each crop, runs the segmentation model and accumulates its
        prediction into a (CPU, padded) buffer of shape $(1, K, H, W)$.
  \item Maintains a count map of how many crops have contributed to each
        pixel.
  \item Returns the buffer divided by the count map, i.e.\ the average
        prediction per pixel.
\end{enumerate}
At training time, the loss is computed on $512{\times}512$ random crops, so
the model never actually sees the full image; the sliding-window protocol at
eval time is what converts that into a coherent full-image prediction. The
overlap of $\tfrac{2}{3}$ is generous and reduces edge artefacts.

For our purposes the key fact is that
\texttt{eval/segmentation/inference.py:make\_inference} wraps
\texttt{slide\_inference} with: an optional rescale-to-target step, an
optional horizontal-flip TTA, and post-processing for the M2F head's
mask-classification output. Our custom-image script (Section~\ref{sec:script})
calls into \texttt{make\_inference} directly.

\subsection{Loss}

\begin{srcbox}\texttt{dinov3/eval/segmentation/loss.py}.\end{srcbox}

\texttt{loss.py} provides two pieces: a weighted-cross-entropy loss with
ignore index, and a soft-Dice loss (mostly disabled in the released
ADE20K config but available). Both follow the convention of returning a
per-element loss that can be combined with a weight tensor. The pattern
\texttt{weighted\_loss} (a decorator factory) wraps a raw element-wise
loss function with weight, reduction, and avg-factor handling. In the
ADE20K linear config only the cross-entropy is active.

% =============================================================================
\section{The Mask2Former head}
\label{sec:m2f}

\begin{srcbox}
\texttt{dinov3/eval/segmentation/models/heads/mask2former\_head.py},\\
\texttt{.../heads/mask2former\_transformer\_decoder.py},\\
\texttt{.../heads/pixel\_decoder.py},\\
\texttt{.../models/utils/} (deformable attention, position encoding, etc.),\\
\texttt{.../models/backbone/dinov3\_adapter.py} (next section),\\
\texttt{configs/config-ade20k-m2f-inference.yaml}.
\end{srcbox}

Mask2Former (Cheng et al.) reformulates segmentation as mask
classification. Instead of predicting a class per pixel, the model
predicts a fixed set of $N$ object queries; each query produces a binary
mask plus a class label. The final per-pixel class is obtained by summing
$\text{mask}_i \cdot p(\text{class}|i)$ across queries. This unifies
semantic, instance, and panoptic segmentation under one architecture and
typically gives substantially higher mIoU than per-pixel classification.

\subsection{Top-level wiring}

The class \texttt{Mask2FormerHead} is constructed inside
\texttt{build\_segmentation\_decoder()} when \texttt{decoder\_type="m2f"}.
It expects an \texttt{input\_shape} dict describing the four feature levels
produced by the DINOv3-Adapter (Section~\ref{sec:adapter}); for ViT-L the
keys are \texttt{"1"}-\texttt{"4"} and the spatial strides decrease from
$P/2$ at the deepest level to $4P$ at the shallowest, where $P$ is the
patch size.

The two main sub-modules are:
\begin{description}[leftmargin=1.6em,style=nextline]
  \item[Pixel decoder.] Implemented in \texttt{pixel\_decoder.py}. It is a
        FPN-with-deformable-attention: the multi-scale features are first
        pushed through a small Transformer encoder that uses
        \emph{multi-scale deformable attention} (MSDeformAttn), then
        progressively upsampled and laterally combined to produce a
        single high-resolution mask feature map (typically at stride 4).
  \item[Transformer decoder.] Implemented in
        \texttt{mask2former\_transformer\_decoder.py}. It takes a fixed
        set of learnable query embeddings and refines them through several
        layers of \emph{masked cross-attention} (each query attends only
        to the predicted-mask region of the encoder features),
        self-attention, and FFN. Each layer outputs intermediate masks +
        class logits used for deep supervision (during training) and
        layer-wise inference.
\end{description}

\subsection{Multi-scale deformable attention}

\begin{srcbox}
\texttt{models/utils/ms\_deform\_attn.py},
\texttt{models/utils/ops/} (the CUDA extension).
\end{srcbox}

Standard cross-attention on dense image features at multiple scales is
quadratic in the number of pixels. Deformable attention replaces it with a
sampling operation: each query predicts a small set of (offset, attention
weight) pairs; the value vectors at those offsets are bilinearly sampled
and combined. The result is linear in the number of queries and effectively
constant per scale. The CUDA op \texttt{MSDeformAttn} implements this
sampling efficiently. Compiling this extension is the single most fragile
step of installing M2F (see Section~\ref{sec:install}).

\subsection{Inference-time output}

The decoder produces, at the last layer, a tensor of \emph{mask logits}
of shape $(B, N, h, w)$ (per-query masks, where $N$ is the number of
queries) and a tensor of \emph{class logits} of shape $(B, N, K{+}1)$
(per-query class scores; the extra slot is a ``no-object'' label). The
inference recipe in \texttt{inference.py:make\_inference} is:
\begin{enumerate}[leftmargin=1.6em,itemsep=0pt]
  \item Softmax the class logits over $K{+}1$, drop the last entry (the
        no-object slot).
  \item Sigmoid the mask logits.
  \item Combine via the einsum
        $\text{out}_{b,c,h,w} = \sum_q p_{b,q,c} \cdot m_{b,q,h,w}$,
        which gives a per-pixel class score map of shape $(B,K,h,w)$.
  \item Argmax over $K$ to get the per-pixel class.
\end{enumerate}
The released ADE20K M2F head also uses test-time augmentation: scales
$[0.9, 0.95, 1.0, 1.05, 1.1]$ averaged, horizontal flips averaged,
sliding-window crops $896{\times}896$ with stride $596$. This is what
gets the model to 55.9 mIoU on ADE20K val (per Meta's model card).

\subsection{What is and isn't shipped}

The M2F training-time machinery (Hungarian matcher, per-mask BCE+Dice
criterion, deep-supervision losses) is present in
\texttt{mask2former\_head.py} for completeness, but the only released
\emph{config} is the inference one. Meta releases a pretrained M2F head
checkpoint (loadable via \texttt{dinov3.hub.segmentors.dinov3\_vit7b16\_ms})
and the recipe to evaluate it. There is no shipped training config for
M2F, so re-training it on a different dataset would require writing one
from scratch. For this project we use M2F as inference-only.

% =============================================================================
\section{The DINOv3-Adapter}
\label{sec:adapter}

\begin{srcbox}
\texttt{dinov3/eval/segmentation/models/backbone/dinov3\_adapter.py}.
\end{srcbox}

This is the part of the M2F path that is novel to DINOv3 (the rest of M2F
is essentially the original Cheng-et-al.\ recipe). Mask2Former needs a
multi-scale feature pyramid; DINOv3 only produces single-scale patch
features (one feature per $16{\times}16$ patch). The adapter manufactures
the missing scales.

\subsection{Approach in plain terms}

Given the ViT's intermediate features at four equally-spaced blocks (the
\texttt{FOUR\_EVEN\_INTERVALS} setting), the adapter:
\begin{enumerate}[leftmargin=1.6em,itemsep=0pt]
  \item Up-samples the shallowest extracted block twice with deconvolutions
        (each $\times 2$) to produce a feature map at stride $P/4$
        ($= 4$ for $P{=}16$). This becomes feature level \texttt{"1"}.
  \item Up-samples the next block once to produce a stride-$P/2$
        ($= 8$) feature map. This becomes level \texttt{"2"}.
  \item Keeps the third extracted block unchanged at stride $P$
        ($= 16$). This is level \texttt{"3"}.
  \item Down-samples the deepest extracted block once with a strided
        convolution to produce a stride-$2P$ ($= 32$) feature map. This
        is level \texttt{"4"}.
\end{enumerate}
The four resulting maps form a feature pyramid spanning strides $4, 8,
16, 32$, exactly what M2F's pixel decoder expects. All four maps share
the same channel dimension as the ViT's embed dim (e.g. 1024 for ViT-L,
4096 for ViT-7B), so the adapter introduces only a few small
deconv/conv layers in addition to the frozen backbone.

\subsection{Why this works at all}

Treating different layers of the ViT as different ``scales'' is a trick
that has been shown to work because deeper ViT layers have larger
effective receptive fields and more semantic features, while shallower
layers retain more local detail. Combined with deconvolutional
up-sampling for the shallow levels and strided convolution for the deep
level, the adapter approximates the multi-scale behaviour of a
hierarchical convnet. The fact that DINOv3 features are already
spatially coherent (thanks to the Gram-anchoring stage during
pretraining) means the upsampled maps actually retain useful local
structure rather than smoothing into uniform blobs. This is the
contribution of the DINOv3-Adapter: a minimal way to inject the
multi-scale assumption back in.

\subsection{Construction details}

\texttt{DINOv3\_Adapter} takes the backbone module and a list
\texttt{interaction\_indexes} (the four block indices to extract). Its
forward pass calls into the backbone's
\texttt{get\_intermediate\_layers(reshape=True)} once, getting a tuple of
four 2D maps, then dispatches each through the appropriate
deconv/identity/conv head and returns a dict keyed by level. This dict is
fed to the M2F pixel decoder.

% =============================================================================
\section{Training the linear head: end-to-end flow}

\begin{srcbox}
\texttt{eval/segmentation/run.py},
\texttt{eval/segmentation/train.py},
\texttt{eval/segmentation/transforms.py},
\texttt{data/datasets/ade20k.py}.
\end{srcbox}

A complete training run, from input to backward pass, is:

\begin{enumerate}[leftmargin=1.6em,itemsep=0pt]
  \item \textbf{Dataset.} \texttt{ADE20K} (in
        \texttt{dinov3/data/datasets/ade20k.py}) reads
        \texttt{ADEChallengeData2016/} via \texttt{torchvision} primitives.
        Each sample is a \texttt{(PIL image, PIL mask)} pair; the mask uses
        ADE20K's index encoding (0 = unlabelled, 1\ldots150 = classes).
  \item \textbf{Paired transforms.} \texttt{transforms.py} defines a
        \texttt{Compose}-style class where every operation acts on
        \texttt{(image, mask)} together: \texttt{RandomResize} (scale ratio
        from a $[0.5, 2.0]$ uniform), \texttt{RandomCrop} to
        $512{\times}512$, \texttt{RandomHorizontalFlip} with $p{=}0.5$, then
        \texttt{Normalize} (image only) with the LVD-1689M mean/std.
  \item \textbf{Forward.}
        \texttt{model.forward()} \emph{is} the wrapped
        \texttt{FeatureDecoder} chain: backbone (frozen, in inference\_mode)
        $\to$ extract intermediate layers $\to$ \texttt{LinearHead.forward()}
        $\to$ raw $(B, K, h, w)$ logits. The training loop passes those
        logits and the resized ground-truth mask to the loss.
  \item \textbf{Loss.} Cross-entropy with \texttt{ignore\_index=255} (the
        value the dataloader uses to mark ``ignore'').
        \texttt{reduce\_zero\_label} in the config makes the dataloader
        subtract 1 from every valid class so that 0 becomes 255 (ignore)
        and class IDs are now $0\ldots149$.
  \item \textbf{Backward.} Standard \texttt{loss.backward()} \texttt{->}
        \texttt{optimizer.step()}. Because the backbone is frozen,
        gradients flow only into the linear head's BatchNorm and
        $1{\times}1$ Conv. This is the reason the head trains so quickly
        despite the large model.
  \item \textbf{Validation.} Every 5000 iterations the eval routine runs
        sliding-window inference on the ADE20K validation set
        (Section~\ref{ssec:slide}) and reports mIoU. The metric is
        accumulated across the full val set in a confusion matrix
        (see \texttt{metrics.py}).
\end{enumerate}

% =============================================================================
\section{Reduced single-GPU recipe}

The released config assumes 8 GPUs. We supply
\texttt{configs/config-ade20k-linear-training-1gpu.yaml} as a single-GPU
variant. We changed \emph{only} hardware-related fields:
\begin{itemize}[leftmargin=1.6em,itemsep=0pt]
  \item \texttt{n\_gpus: 1}, \texttt{bs: 1}: effective batch 1 instead of
        16. With a frozen backbone and a tiny head this is OK in practice
        (BatchNorm sees only one sample per step, but the head's
        BatchNorm is a single-channel per-class scaling that can be
        replaced with a no-op or kept as-is for evaluation).
  \item \texttt{transforms.train.img\_size: 512}: same as the multi-GPU
        config; can be reduced to 448 if the chosen backbone is too
        large for the GPU.
  \item Crop size and stride at eval are unchanged (512 / 341).
\end{itemize}
Algorithmic settings -- optimizer, scheduler, loss weights, head structure,
\texttt{reduce\_zero\_label}, etc.\ -- are preserved exactly. Whether the
1-GPU run actually converges to mIoU within a few points of the 8-GPU run
depends on the smaller effective batch size; for our class demo we are
comfortable with the qualitative result and the partial-training
checkpoint as a teaching artefact.

% =============================================================================
\section{Custom-image inference}
\label{sec:script}

The repository accompanying this report contains one script that does not
exist in Meta's repo: \texttt{scripts/segment\_image.py}. Its only role is
to invoke the existing pieces correctly on a single user image.

The pipeline it implements:
\begin{enumerate}[leftmargin=1.6em,itemsep=0pt]
  \item Load the image with PIL, normalise with the LVD-1689M
        mean/std (since the released backbones are pretrained on web
        images), optionally resize the short side to a target size.
  \item Build the chosen backbone via Meta's hub factory
        (\texttt{dinov3.hub.backbones.dinov3\_vits16}, etc.), passing the
        local checkpoint path as \texttt{weights=}.
  \item Build the chosen decoder via
        \texttt{dinov3.eval.segmentation.models.build\_segmentation\_decoder}
        with \texttt{decoder\_type=\{linear,m2f\}}; the function returns
        a \texttt{FeatureDecoder} that already wraps the backbone + head.
  \item Load the decoder checkpoint with \texttt{strict=False}, mirroring
        the way Meta's \texttt{dinov3/hub/segmentors.py} loads the
        released M2F head (the backbone keys are skipped because the
        backbone was loaded separately).
  \item Call \texttt{dinov3.eval.segmentation.inference.make\_inference}
        with \texttt{inference\_mode="slide"} and the head-appropriate
        crop / stride: $(512, 341)$ for the linear head and $(896, 596)$
        for M2F, matching Meta's eval configs.
  \item Convert the output logits to a class map via \texttt{argmax},
        colour-code the classes, and overlay on the original image.
\end{enumerate}
The script is roughly 300 lines, of which the modelling-relevant parts
are simple wrappers around the four \texttt{import}s above. The colour
palette and the overlay blending are routine numpy. We deliberately do
not bundle ADE20K's official palette file with the repository; instead
we generate a deterministic palette via HSV with a fixed seed so the
visualisations are reproducible.

% =============================================================================
\section{Installing it: the practical bits}
\label{sec:install}

The project's \texttt{README.md} has the full walkthrough. Key practical
points:

\paragraph{Cloning vs.\ vendoring.} We \emph{do not} include Meta's source
in this repository. Instead, the user clones
\texttt{facebookresearch/dinov3} alongside, installs it with
\texttt{pip install -e .}, and our scripts import it as a normal package.
This keeps the licence story clear and the repository small.

\paragraph{The MSDeformAttn extension.} Mask2Former requires a custom CUDA
op compiled at install time (\texttt{python setup.py build\_ext --inplace}
in \texttt{dinov3/eval/segmentation/models/utils/ops/}). The most common
failure modes are: (a) CUDA toolkit version mismatched with the PyTorch
build, (b) \texttt{nvcc} not on \texttt{PATH}, (c)
\texttt{TORCH\_CUDA\_ARCH\_LIST} unset on machines whose driver
mis-reports the device capability. The op is needed only for M2F; the
linear head runs without it.

\paragraph{Pretrained weights.} Three to four checkpoints in total: a
backbone (any of S/B/L/7B), the released ADE20K M2F head, and any linear
head that you train yourself. The backbones are gated behind a Meta
download form. The data (ADE20K) is gated behind the MIT release form.
Neither is redistributable.

% =============================================================================
\section{Expected results}

The headline numbers, taken from Meta's \texttt{MODEL\_CARD.md}, are:
\begin{center}\small
\begin{tabular}{lcc}
\toprule
Model & ADE20K mIoU & Notes \\
\midrule
DINOv3 ViT-L/16 + linear head      & 54.9 & frozen backbone, $1{\times}1$ Conv head \\
DINOv3 ViT-7B/16 + linear head     & 55.9 & frozen backbone, $1{\times}1$ Conv head \\
DINOv3 ViT-7B/16 + M2F head        & 55.9 & released checkpoint, TTA at 5 scales \\
\bottomrule
\end{tabular}
\end{center}
That the 7B linear head matches the 7B M2F head on this metric is, by
itself, a striking result -- it suggests that for ADE20K the bottleneck is
the backbone, not the decoder. M2F retains its advantage on instance- and
panoptic-segmentation metrics that the linear head cannot reach by
construction (the linear head always assigns the same class to all pixels
in a region; it has no notion of separate instances).

For our class project we run two demonstrations:
\begin{enumerate}[leftmargin=1.6em,itemsep=0pt]
  \item Train the linear head on ADE20K (single-GPU recipe), report
        train/val curves, and evaluate at the iteration we can afford to
        reach (typically a few thousand iterations on a single consumer
        GPU). The checkpoint we obtain will not match the 54.9 mIoU figure
        above but should clearly demonstrate the training behaviour and
        produce sensible qualitative outputs on ADE20K val and on
        custom images.
  \item Run M2F inference at full quality on ADE20K val and on custom
        images, reproducing the qualitative results from the model card.
\end{enumerate}

% =============================================================================
\section{Why DINOv3 features are well-suited to segmentation}

Three properties of the DINOv3 backbone matter directly for segmentation
quality:

\paragraph{Patch features are spatially coherent.} The Gram-anchoring
training stage explicitly preserves the pairwise patch-similarity
structure across long training, which is the property a segmentation head
exploits when grouping nearby patches into the same class.

\paragraph{Register tokens absorb high-norm artefacts.} In ViTs trained
without registers, a small number of patch tokens become high-norm
``sinks'' that hold global information; those positions appear as bright
artefacts in segmentation outputs. DINOv3 includes 4 register tokens
explicitly to absorb this role, leaving patch tokens spatially clean and
class-discriminative.

\paragraph{High-resolution adaptation.} The third pretraining stage trains
the model at multiple resolutions including $768$, so applying it at
ADE20K's typical $512$--$896$ inputs is in-distribution. With learnable
absolute positional embeddings (DINOv2) this would have required
interpolation; with axial RoPE (DINOv3) it is native.

These three together mean that a $1{\times}1$ convolution on top of frozen
DINOv3 features achieves a result that, three years ago, would have
required a heavy task-specific decoder.

% =============================================================================
\section{File-by-file map}

\begin{center}\footnotesize
\begin{tabular}{ll}
\toprule
File in \texttt{dinov3/} & Role \\
\midrule
\texttt{models/vision\_transformer.py}                 & ViT backbone with the \texttt{get\_intermediate\_layers} hook \\
\texttt{eval/utils.py}                                 & \texttt{ModelWithIntermediateLayers} (frozen-backbone wrapper) \\
\texttt{eval/segmentation/models/\_\_init\_\_.py}      & dispatcher: linear vs M2F \\
\texttt{eval/segmentation/models/heads/linear\_head.py}& linear segmentation head \\
\texttt{eval/segmentation/models/heads/mask2former\_head.py} & M2F top-level head \\
\texttt{eval/segmentation/models/heads/mask2former\_transformer\_decoder.py} & M2F transformer decoder \\
\texttt{eval/segmentation/models/heads/pixel\_decoder.py} & M2F pixel decoder (FPN + MSDeformAttn) \\
\texttt{eval/segmentation/models/backbone/dinov3\_adapter.py} & multi-scale adapter for M2F \\
\texttt{eval/segmentation/models/utils/ms\_deform\_attn.py} & deformable attention layer \\
\texttt{eval/segmentation/models/utils/ops/}           & CUDA extension for MSDeformAttn \\
\texttt{eval/segmentation/inference.py}                & sliding-window + TTA inference \\
\texttt{eval/segmentation/train.py, eval.py, run.py}   & training and evaluation loops \\
\texttt{eval/segmentation/loss.py, metrics.py}         & cross-entropy/Dice losses, mIoU \\
\texttt{eval/segmentation/transforms.py}               & paired (image, mask) augmentations \\
\texttt{eval/segmentation/configs/}                    & ADE20K linear-training and M2F-inference YAMLs \\
\texttt{hub/segmentors.py}                             & prebuilt M2F segmentor with weight downloading \\
\bottomrule
\end{tabular}
\end{center}

\bigskip
\noindent
This concludes the report. The companion repository contains the README,
the single-GPU config, and the custom-image inference script that wires
all of the above into a runnable demo on user-supplied images.

\end{document}
